\section{Measure-theoretic details and extended-value identities}\label{app:measure}
\subsection{Exactly which foundational results are used}
Standard Borel disintegration, measurable integration of kernels, and countable sequential-kernel extension are used as foundational mathematical results \cite{Kallenberg}. They assert existence and almost-everywhere uniqueness of a conditional kernel, not a common differentiable version across all parameters. In the continuous-alphabet theorem such versions are provided by positive continuous densities; in the geometric theorem derivatives are constructed as currents instead.

For a measurable report $R$, a disintegration $P_y$ satisfies
\[
 P(A\cap R^{-1}(B))=\int_B P_y(A)P_R(dy).
\]
If $Q=qP$, then $Q_R=bP_R$ with $b=E(q\mid R)$, and on $b>0$, the conditional density of $Q_y$ relative to $P_y$ is $q/b$. To verify it, integrate against rectangle tests $\one_A\one_B(R)$ and use the definition of conditional expectation. When $b=0$, $q=0$ on that fiber almost surely, so it contributes nothing to $Q$.

\subsection{Entropy as a monotone finite-partition limit}
Let $\mathcal A_m$ be increasing finite sigma-fields generating the relevant countably generated sigma-field modulo completion. For $q=dQ/dP$ set $q_m=E_P(q\mid\mathcal A_m)$. With $\phi(t)=t\log t-t+1\ge0$, Jensen shows that $E\phi(q_m)$ is increasing and at most $E\phi(q)$. Conditional martingale convergence gives $q_m\to q$ almost surely and in $L^1$. Fatou supplies
$E\phi(q)\le\liminf_mE\phi(q_m)$, hence
\[
 D(Q\Vert P)=\lim_mD(Q|_{\mathcal A_m}\Vert P|_{\mathcal A_m}),
\]
including $+\infty$. The entropy chain rule can be checked either through this approximation or through the integrable-negative-part argument in Theorem~\ref{thm:entropy-chain}; neither proof subtracts infinite entropies.

\begin{lemma}[Pinsker's bound in the convention of this paper]\label{lem:pinsker}
For probabilities $P,Q$,
\[
 d_{\TV}(Q,P)^2\le\tfrac12D(Q\Vert P).
\]
\end{lemma}
\begin{proof}
For an event $A$, conditional Jensen bounds the full entropy below by the Bernoulli entropy $d(q\Vert p)$ with $q=Q(A),p=P(A)$. For $0<p<1$,
$\partial_q^2d(q\Vert p)=1/[q(1-q)]\ge4$, and both $d$ and its first derivative vanish at $q=p$. Integrating twice gives $d(q\Vert p)\ge2(q-p)^2$. Boundary cases follow by continuity or infinite entropy. Take the supremum over $A$.
\end{proof}

\subsection{Conditional source derivatives and common versions}
A bounded source satisfies
$e^{\inf V}\le E(e^V\mid R)\le e^{\sup V}$ almost surely. It therefore has a bounded logarithmic conditional version. For finitely many bounded sources, conditional exponential series can be dominated on every compact source window. Countably many rational coefficients and the resulting continuous extensions furnish a common version for their finite source derivatives. This assertion does not extend to arbitrary physical parameters or parameter-dependent reports without the separately proved transport or density construction.

For a positive numerator/denominator pair, the derivative of a conditional expectation is a quotient derivative. Multiplying back by the evidence gives an integrated estimate. An arbitrary out-of-range version on a zero-evidence set can be nonzero under a second model and invalidate cross-model comparisons. Throughout the paper bounded payoffs use bounded posterior versions, with zero an admissible choice.

\subsection{A direct proof of the experimental test chain rule}
For prescribed actions $u_{0:m-1}$, the observation density is the product of its successive conditional densities. Under a policy, the joint action--observation measure is this observation product multiplied by the chronological policy kernels. This follows by repeated disintegration of \eqref{eq:instrument-path}. In a finite alphabet, summing over action strings yields the adaptive observation law. In the continuous theorem, the same product is a jointly measurable positive density; Tonelli allows its integration against policy kernels. The density is not the conditional law given the \emph{final} adaptive design. For example, when the next action equals the previous observation, conditioning on that action already reveals the earlier readout.

\subsection{Source truncation when a bounded-source theorem is used}
For an unbounded $V$ with $0<Pe^V<\infty$, clipped sources $V_m=(-m)\vee(V\wedge m)$ satisfy $e^{V_m}\to e^V$ and are dominated by $1+e^V$. Thus their normalizations and selected laws converge in total variation by dominated convergence and Lemma~\ref{lem:normalization}. Derivatives require domination of the derivative-weighted exponentials as well. A finite exponential moment at one parameter does not provide an entire neighborhood of differentiability.
\par
